% Section 5.1, from lectures/L05/appendix/a_device_drivers.md.

\section{Device Drivers}
\label{c5:sec:devicedrivers}\label{c5:app:a}
This section covers how Zephyr turns a devicetree node into a device, what a driver provides for it,
and how a driver for a new part is packaged: binding, Kconfig option, and module. By the end you
should be able to explain every file in \file{lectures/L05/examples/thermo/module}, and why each one
is needed.

\subsection{What a driver is, in Zephyr}
Chapter~\ref{c2:ch} used devices without asking where they came from:
\code{DEVICE_DT_GET(DT_ALIAS(led0))} returned a \code{const struct device *}, and the GPIO API did
the rest. This chapter is the other side of that call.

A Zephyr driver is three things:

\begin{enumerate}
  \item \textbf{A set of functions} that implement one of Zephyr’s device APIs, such as GPIO, I2C or sensor,
    for one kind of part.
  \item \textbf{An init function}, which the kernel calls once at boot to bring an instance of the part up.
  \item \textbf{A macro expansion per devicetree node}, which creates a \code{struct device} for every node
    that the driver serves, at compile time, with the node’s properties already filled in.
\end{enumerate}

There is no registration at run time and no probing. By the time the program links, every device
exists as a constant in flash, and the only thing left to happen at boot is the init function.

\subsection{\code{struct device}}
Every device is a \code{struct device}, defined in \code{<zephyr/device.h>}. Its first five members
are the ones a driver writer deals with:

\begin{ccode}
struct device {
	const char *name;           /* The devicetree node's name, e.g. "thermo@48". */
	const void *config;         /* Per-instance config: const, from devicetree. */
	const void *api;            /* The function table: the driver's API struct. */
	struct device_state *state; /* Owned by the kernel: init's state and result. */
	void *data;                 /* Per-instance run-time data: mutable. */
	/* ... */
};
\end{ccode}

\noindent \textbf{\code{config} and \code{data} are the driver’s own types}, which the kernel stores
as \code{void *}. The difference between them matters on a microcontroller:

\begin{itemize}
  \item \textbf{\code{config} is \code{const}.} It holds what the devicetree says about the instance: its
    bus, its address, its properties. It is built at compile time and lives in flash.
  \item \textbf{\code{data} is not.} It holds what changes at run time: the last reading, a mutex, a state
    machine. It lives in RAM.
\end{itemize}

In \file{lectures/L05/examples/thermo/module/drivers/sensor/thermo/thermo.c} they are:

\begin{ccode}
struct thermo_config {
	struct i2c_dt_spec i2c;
	uint32_t conversion_timeout_ms;
};

struct thermo_data {
	int16_t raw;
};
\end{ccode}

\noindent and every driver function starts by getting them back out of the device:

\begin{ccode}
const struct thermo_config *config = dev->config;
struct thermo_data *data = dev->data;
\end{ccode}

\noindent \textbf{The name is the node’s name}, \code{thermo@48}, not its label. That is what
\code{dev->name} prints and what the logs show.

\subsection{The API struct}
Each kind of device has an API: a struct of function pointers declared by its subsystem. For the
sensor subsystem it is \code{struct sensor_driver_api} in \code{<zephyr/drivers/sensor.h>}, whose
two mandatory members are:

\begin{ccode}
sensor_sample_fetch_t sample_fetch;
sensor_channel_get_t channel_get;
\end{ccode}

\noindent The functions an application calls, such as \code{sensor_sample_fetch(dev)}, are thin
wrappers that cast \code{dev->api} and call through it:

\begin{ccode}
static inline int z_impl_sensor_sample_fetch(const struct device *dev)
{
	const struct sensor_driver_api *api =
		(const struct sensor_driver_api *)dev->api;

	return api->sample_fetch(dev, SENSOR_CHAN_ALL);
}
\end{ccode}

A driver fills in its own table with \code{DEVICE_API()}:

\begin{ccode}
static DEVICE_API(sensor, thermo_api) = {
	.sample_fetch = thermo_sample_fetch,
	.channel_get = thermo_channel_get,
};
\end{ccode}

\noindent \code{DEVICE_API(sensor, name)} declares a \code{const struct sensor_driver_api}, and also
places it in a linker section of its own, one per API class. That lets \code{DEVICE_API_IS(sensor,
dev)} check at run time whether a device really has a sensor API, rather than trusting whoever
passed the pointer.

This is Zephyr’s polymorphism: the application’s \code{sensor_sample_fetch()} works the same for
Thermo, for a Bosch accelerometer, or for any of the hundreds of sensor drivers in the tree, and
none of them knows about the others.

\subsection{From devicetree node to device}

\subsubsection{\code{DT_DRV_COMPAT} and instances}
A driver serves every node with a given \code{compatible}, and says which one at the top of its
file:

\begin{ccode}
#define DT_DRV_COMPAT qacademy_thermo
\end{ccode}

\noindent The value is the compatible \code{"qacademy,thermo"} in lowercase, with every character
that cannot be in a C identifier turned into an underscore. From then on, the \code{DT_INST_...}
macros count the \textbf{instances} of that compatible whose \code{status} is \code{okay}, numbered
0, 1, 2, … in an order the build decides. \code{DT_INST_PROP(inst, conversion_timeout_ms)} is
\code{DT_PROP()} for instance \code{inst}.

\subsubsection{One device per instance}
The driver writes a macro that defines everything for one instance, and expands it for all of them:

\begin{ccode}
#define THERMO_DEFINE(inst)                                                 \
    static struct thermo_data thermo_data_##inst;                           \
                                                                            \
    static const struct thermo_config thermo_config_##inst = {              \
        .i2c = I2C_DT_SPEC_INST_GET(inst),                                  \
        .conversion_timeout_ms = DT_INST_PROP(inst, conversion_timeout_ms), \
    };                                                                      \
                                                                            \
    SENSOR_DEVICE_DT_INST_DEFINE(inst, thermo_init, NULL,                   \
                                 &thermo_data_##inst,                       \
                                 &thermo_config_##inst, POST_KERNEL,        \
                                 CONFIG_SENSOR_INIT_PRIORITY, &thermo_api);

DT_INST_FOREACH_STATUS_OKAY(THERMO_DEFINE)
\end{ccode}

\noindent With two Thermo nodes in the devicetree, that is two \code{thermo_data}, two
\code{thermo_config}, and two devices; with none, it is nothing at all.
\code{I2C_DT_SPEC_INST_GET(inst)} fills in the bus device and the address from the node’s parent and
its \code{reg} property.

The device macro takes, in order:

\begin{booktable}{@{}LL@{}}
  \thead{Argument} & \thead{Meaning} \\
  \midrule
  \code{inst} & The instance. The macro turns it into the node, and the node into the device’s name and symbol. \\
  \code{thermo_init} & The init function. \\
  \code{NULL} & Power management actions, which Thermo has none of. \\
  \code{&thermo_data_\#\#inst} & \code{dev->data}. May be \code{NULL}. \\
  \code{&thermo_config_\#\#inst} & \code{dev->config}. \\
  \code{POST_KERNEL} & The init level. \\
  \code{CONFIG_SENSOR_INIT_PRIORITY} & The priority within the level. \\
  \code{&thermo_api} & \code{dev->api}. \\
\end{booktable}

\code{SENSOR_DEVICE_DT_INST_DEFINE()} is \code{DEVICE_DT_INST_DEFINE()} plus the sensor subsystem’s
bookkeeping; other subsystems’ drivers use \code{DEVICE_DT_INST_DEFINE()} directly.

\subsubsection{The symbol behind \code{DEVICE_DT_GET}}
A device is a global variable named after the node’s position in the devicetree’s dependency order:
\code{__device_dts_ord_32} is the device of the node with ordinal 32, and the ordinals are listed at
the top of \code{build/zephyr/include/generated/zephyr/devicetree_generated.h}.
\code{DEVICE_DT_GET(node)} is a reference to that variable.

That explains the error Chapter~\ref{c2:ch} met, from the other side. \code{DEVICE_DT_GET()} of an
enabled node that no driver defined a device for compiles, since the name is only a symbol, but the
link fails:

\begin{console}
undefined reference to `__device_dts_ord_32'
\end{console}

\noindent The node exists; nothing defined its device. The causes are a driver that is not enabled,
a driver whose \code{DT_DRV_COMPAT} does not match the node, or a module that is not in the build.
In the Thermo example, removing the module from \code{CMakeLists.txt} does exactly that.

\subsection{Initialization}

\subsubsection{Levels and priorities}
The kernel initializes devices in levels, in this order:

\begin{booktable}{@{}LL@{}}
  \thead{Level} & \thead{When} \\
  \midrule
  \code{EARLY} & Right after entering C, before the kernel’s architecture-specific initialization. \\
  \code{PRE_KERNEL_1}, \code{PRE_KERNEL_2} & Before the kernel’s services are available: no sleeping, no kernel objects. Interrupt controllers, clocks, the system timer, consoles. \\
  \code{POST_KERNEL} & After the kernel is running. Most device drivers, since they may need to sleep or take a mutex. \\
  \code{APPLICATION} & Just before \code{main()}. \\
\end{booktable}

Within a level, devices run in order of their \textbf{priority}, a number from 0 to 999, lowest
first. Subsystems give their drivers a Kconfig option for it: I2C controllers use
\code{CONFIG_I2C_INIT_PRIORITY}, which defaults to \code{CONFIG_KERNEL_INIT_PRIORITY_DEVICE} (50),
and sensors use \code{CONFIG_SENSOR_INIT_PRIORITY} (90). A sensor on an I2C bus therefore
initializes after the bus, which it needs.

\textbf{The build checks that.} \code{CONFIG_CHECK_INIT_PRIORITIES}, on by default, compares the
init order in the image with the dependencies in the devicetree, and fails the build if a device
initializes before a device its node depends on, such as a sensor before its bus.

To see the order, list the build’s init entries:

\begin{shell}
west build -t initlevels
\end{shell}

\noindent For the Chapter~\ref{c5:ch} exercises’ test application with two Lumo parts, the end of
the list is:

\begin{console}
POST_KERNEL
  __init_enable_logger: enable_logger(NULL)
  __init___device_dts_ord_31: i2c_emul_init(__device_dts_ord_31)
  __init___device_dts_ord_32: lumo_init(__device_dts_ord_32)
  __init___device_dts_ord_33: lumo_init(__device_dts_ord_33)
\end{console}

\noindent The list is read from the image, so it also shows a device whose initialization is
deferred (below), which the kernel then skips at boot.

\subsubsection{The init function}
An init function takes the device and returns 0 or a negative error number:

\begin{ccode}
static int thermo_init(const struct device *dev)
{
	const struct thermo_config *config = dev->config;
	uint8_t id;
	int ret;

	if (!i2c_is_ready_dt(&config->i2c)) {
		return -ENODEV;
	}

	ret = i2c_reg_read_byte_dt(&config->i2c, THERMO_REG_WHO_AM_I, &id);
	if (ret < 0) {
		return ret;
	}

	if (id != THERMO_WHO_AM_I_VALUE) {
		return -ENODEV;
	}

	return 0;
}
\end{ccode}

\noindent What belongs in it: checking that the bus is ready, that the part answers, and that it is
the part the devicetree says it is; and putting the part into the state the devicetree asks for.
What does not: anything slow that a driver could do on first use instead. Every init function runs
before \code{main()}, one after another.

\subsubsection{\code{device_is_ready()}}
The kernel records whether a device’s init function ran, and what it returned.
\code{device_is_ready(dev)} is true only if both went well:

\begin{ccode}
return dev->state->initialized && (dev->state->init_res == 0U);
\end{ccode}

\noindent \textbf{A failed init is final.} Nothing retries it, and every call into the device’s API
would talk to a part that is not there or not right. That is why every user of a device checks
\code{device_is_ready()} once, before using it, and why a driver’s init must fail when the part is
wrong rather than log a warning and carry on.

\subsubsection{Deferred initialization}
A node with the property \code{zephyr,deferred-init} is not initialized at boot. The application
initializes it later, when it chooses, with \code{device_init(dev)}: for example once a power rail
has been switched on. \code{device_init()} returns what the init function returned, or
\code{-EALREADY} if the device was initialized before, \textbf{whether or not that went well}. The
Thermo tests use exactly this to test a failing init: they break the emulated part, call
\code{device_init()}, and get \code{-ENODEV}; a second call, after repairing the part, gets
\code{-EALREADY}, and the device stays not ready.

\subsection{A binding and a Kconfig option for a new part}

\subsubsection{The binding}
Chapter~\ref{c2:ch} wrote a binding for a node of the application’s own. A driver’s binding is the
same kind of file, and says what a node for the part must look like:

\begin{yamlcode}
description: |
  Thermo, a made-up I2C temperature sensor ...

compatible: "qacademy,thermo"

include: [sensor-device.yaml, i2c-device.yaml]

properties:
  conversion-timeout-ms:
    type: int
    default: 100
\end{yamlcode}

\noindent \code{i2c-device.yaml} is Zephyr’s common binding for a part on an I2C bus: it requires
\code{reg}, the address, and makes the node valid only as a child of an I2C controller.
\code{sensor-device.yaml} adds the properties every sensor may have. A property with an \code{enum:}
list may only take one of the listed values, and \code{DT_INST_ENUM_IDX(inst, prop)} gives the index
of the value the node chose, at compile time.

\subsubsection{The vendor prefix}
The part before the comma in a compatible is a \textbf{vendor prefix}, and Zephyr keeps a list of
known ones in \code{dts/bindings/vendor-prefixes.txt}. A module can add its own, in a file of the
same name:

\begin{console}
qacademy	QRTECH Academy
\end{console}

\noindent Without it, a normal build only warns, once per node:

\begin{console}
node '/i2c@100/lumo@29' compatible 'qacademy,lumo' has unknown vendor prefix 'qacademy'
\end{console}

\noindent \textbf{Twister turns that warning into an error.} It configures every build with
\code{-DEXTRA_GEN_EDT_ARGS=--edtlib-Werror}, and the same message becomes a \code{devicetree error:}
that stops the configuration. A module that builds fine by hand and fails under Twister has usually
missed exactly this file.

\subsubsection{The Kconfig option}
Every driver has a Kconfig option that decides whether it is compiled:

\begin{kconfigcode}
config THERMO
	bool "Thermo I2C temperature sensor"
	default y
	depends on SENSOR
	depends on DT_HAS_QACADEMY_THERMO_ENABLED
	select I2C
\end{kconfigcode}

\noindent \textbf{\code{DT_HAS_QACADEMY_THERMO_ENABLED} is generated from the devicetree.} For every
compatible that has a binding, the build writes a symbol into \code{build/Kconfig/Kconfig.dts} that
is \code{y} when at least one node with that compatible is enabled:

\begin{console}
DT_COMPAT_QACADEMY_LUMO := qacademy,lumo

config DT_HAS_QACADEMY_LUMO_ENABLED
	def_bool $(dt_compat_enabled,$(DT_COMPAT_QACADEMY_LUMO))
\end{console}

\noindent With \code{default y}, the driver is on exactly when the hardware is there, and nobody
sets \code{CONFIG_THERMO} in a \code{prj.conf}. This is the pattern of nearly every driver in
Zephyr, and the reason enabling a driver usually means adding a devicetree node rather than setting
an option. The remaining dependencies are the subsystem the driver plugs into (\code{SENSOR}, which
the application enables) and the bus it needs (\code{I2C}, which the driver selects).

One consequence is easy to miss. A \code{default} that does not take is not an assignment, so
nothing warns about it: an application with a Thermo node but without \code{CONFIG_SENSOR=y} builds
its devicetree, leaves \code{THERMO} off, and fails at link time with \code{undefined reference to
'__device_dts_ord_32'}.

\subsection{The sensor API}

\subsubsection{Fetch and get}
Reading a sensor takes two calls:

\begin{ccode}
ret = sensor_sample_fetch(dev);
ret = sensor_channel_get(dev, SENSOR_CHAN_AMBIENT_TEMP, &value);
\end{ccode}

\noindent \textbf{\code{sample_fetch} talks to the part}: it starts a measurement if the part needs
one, reads the raw result, and stores it in \code{dev->data}. It can fail, and it can take time.
\code{sensor_sample_fetch(dev)} fetches every channel the part has (\code{SENSOR_CHAN_ALL}), and
\code{sensor_sample_fetch_chan(dev, chan)} only one.

\textbf{\code{channel_get} only converts} the stored result into the channel’s unit. It does no I/O,
and it returns the same value every time until the next fetch. The split lets a part that measures
several things at once, such as an accelerometer’s three axes, be read in one bus transaction and
then converted channel by channel. Both return \code{-ENOTSUP} for a channel the part does not have.

\subsubsection{\code{struct sensor_value}}
A reading is a fixed-point number in two parts:

\begin{ccode}
struct sensor_value {
	int32_t val1;   /* Integer part. */
	int32_t val2;   /* Fractional part, in millionths. */
};
\end{ccode}

\noindent Each channel has a unit: degrees Celsius for \code{SENSOR_CHAN_AMBIENT_TEMP}, lux for
\code{SENSOR_CHAN_LIGHT}. \textbf{Both parts carry the sign}: -1.5 degrees is \code{val1 = -1, val2
= -500000}, and -0.5 degrees is \code{val1 = 0, val2 = -500000}. A driver should not assemble the
two parts by hand;
\code{sensor_value_from_micro(&value, micro)} does it from a value in millionths, and
\code{sensor_value_to_micro(&value)} converts back. Work in \code{int64_t} millionths: a 16-bit count times
250000 millionths per count does not fit in 32 bits.

\subsubsection{The newer, asynchronous API}
Zephyr also has an asynchronous sensor API (\code{CONFIG_SENSOR_ASYNC_API}), built on RTIO, in which
a driver provides \code{submit} and \code{get_decoder} instead, and readings are decoded into Q31
fixed-point values. It suits high-rate and streaming sensors. The fetch and get API is what most
existing drivers implement, and what this book uses; Q31 values reappear in \secref{c5:app:b}, in
the emulator backend.

\subsection{Zephyr modules}

\subsubsection{Out of tree}
A driver for a part that Zephyr does not support does not have to live in Zephyr. A \textbf{module}
is a directory that the build treats as an extension of Zephyr: its CMake is included into Zephyr’s
build, its Kconfig into Zephyr’s configuration, and its bindings into the devicetree search path.
Vendor HALs, mbedTLS and LVGL are all modules; so is \code{examples/thermo/module}.

What makes a directory a module is one file, \code{zephyr/module.yml}:

\begin{yamlcode}
name: qacademy-thermo
build:
  cmake: .
  kconfig: Kconfig
  settings:
    dts_root: .
\end{yamlcode}

\begin{booktable}{@{}lL@{}}
  \thead{Key} & \thead{Meaning} \\
  \midrule
  \code{name} & The module’s name. The build also defines a Kconfig symbol for it: \code{CONFIG_ZEPHYR_QACADEMY_THERMO_MODULE=y}. \\
  \code{build: cmake} & The directory whose \code{CMakeLists.txt} the build includes. \\
  \code{build: kconfig} & The Kconfig file the build sources. \\
  \code{settings: dts_root} & A directory whose \code{dts/bindings/} holds bindings, and whose \code{dts/bindings/vendor-prefixes.txt} holds vendor prefixes. \\
\end{booktable}

The module’s \code{CMakeLists.txt} is included after Kconfig has run, so it compiles the driver only
when the driver is enabled:

\begin{cmakecode}
add_subdirectory_ifdef(CONFIG_THERMO drivers/sensor/thermo)
\end{cmakecode}

\noindent and the driver’s own \code{CMakeLists.txt} makes it a library of the Zephyr image, as
in-tree drivers do:

\begin{cmakecode}
zephyr_library()
zephyr_library_sources(thermo.c)
\end{cmakecode}

\subsubsection{Adding a module to a build}
In a west workspace, the modules are the projects of the manifest, as \secref{c1:app:b} described. A
module that is not a west project is added by path, in \code{ZEPHYR_EXTRA_MODULES}, before
\code{find_package(Zephyr)}:

\begin{cmakecode}
list(APPEND ZEPHYR_EXTRA_MODULES ${CMAKE_CURRENT_SOURCE_DIR}/../module)
find_package(Zephyr REQUIRED HINTS $ENV{ZEPHYR_BASE})
\end{cmakecode}

\noindent or on the command line, as \code{-DZEPHYR_EXTRA_MODULES=<path>}. The build merges the two,
so an application can name the modules it always needs in its \code{CMakeLists.txt}, and a test can
add another, such as an emulator, from outside.

\subsection{Summary}
\begin{itemize}
  \item A driver is an API implementation, an init function, and a macro expanded once for every enabled
    devicetree node of its compatible.
  \item A device’s \code{config} is \code{const}, from devicetree, in flash; its \code{data} changes at run
    time, in RAM.
  \item \code{DT_DRV_COMPAT} names the compatible; \code{DT_INST_FOREACH_STATUS_OKAY()} and
    \code{DEVICE_DT_INST_DEFINE()} make one device per instance.
  \item \code{undefined reference to '__device_dts_ord_N'} means a node that no driver made a device for.
  \item Devices initialize by level, then by priority. \code{west build -t initlevels} shows the order, and
    the build checks it against the devicetree.
  \item A failed init is final: \code{device_is_ready()} is false from then on. \code{zephyr,deferred-init}
    and \code{device_init()} move init to a moment the application chooses.
  \item A new part needs a binding, a vendor prefix, and a Kconfig option that depends on
    \code{DT_HAS_<COMPAT>_ENABLED}, so that the devicetree node alone enables the driver.
  \item The sensor API splits I/O (\code{sample_fetch}) from conversion (\code{channel_get}), and returns
    fixed-point \code{sensor_value}s whose two parts both carry the sign.
  \item A module packages drivers, bindings and Kconfig outside Zephyr, and is added to a build with
    \code{ZEPHYR_EXTRA_MODULES} or a west manifest.
\end{itemize}
